% ============================================================
% paper.tex -- draft manuscript (v0.1, 17 November 2026)
% Compile with TinyTeX:  tinytex::latexmk("manuscript/paper.tex")
% Status: METHODS DOCUMENTATION DRAFT. The literature survey is not yet
%         done; every point needing real citations is marked \todo{}.
% ============================================================
\documentclass[11pt,a4paper]{article}
\usepackage[margin=2.4cm]{geometry}
\usepackage[T1]{fontenc}
\usepackage[utf8]{inputenc}
\usepackage{booktabs}
\usepackage{graphicx}
\usepackage{xcolor}
\usepackage[hidelinks]{hyperref}
\usepackage{caption}

\newcommand{\todo}[1]{\textcolor{red}{[TODO: #1]}}

\title{\textbf{A two-stage framework for selecting CMIP6 Earth-system models\\
for catchment-scale hydrological application over Peninsular India}\\
\large Empirical quantile mapping onto CAMELS-IND forcings, and a GR4J-based
discharge evaluation for out-of-sample model shortlisting}
\author{Ankit Deshmukh\thanks{Draft v0.1 -- methods documentation. Objective
statement and literature survey pending (Section~\ref{sec:todo}).}}
\date{17 November 2026}

\begin{document}
\maketitle

\begin{abstract}
\noindent
Earth-system models participating in CMIP6 differ strongly in how
realistically they reproduce the hydro-climate of Peninsular India, yet
scenario-driven catchment water-availability studies typically use model
ensembles selected arbitrarily or inherited from other regions. We develop a
reproducible two-stage model-selection framework over the 472 CAMELS-IND
catchments. Stage one bias-corrects daily CMIP6 precipitation and
temperature (1981--2014, 22 models with all three required variables) onto
catchment-mean forcings with seasonal empirical quantile mapping, and ranks
models on \emph{out-of-sample} forcing skill (fit 1981--2000, score
2001--2014). Stage two drives a lumped hydrological model (GR4J, calibrated
per catchment on observed discharge) with each bias-corrected series and
ranks models by how realistically the resulting \emph{discharge} matches
observed flows. A hybrid of the two rankings yields a defensible 14-model
shortlist spanning the CMIP6 equilibrium-climate-sensitivity range
(1.8--5.7\,K) for future SSP1-2.6/SSP2-4.5/SSP5-8.5 water-availability and
drought analyses. Calibration of GR4J across 192 gauged catchments reaches
a median calibration Kling--Gupta efficiency of 0.80 and out-of-sample
(validation) median of 0.54. We document the honest handling of model
gaps (365-day calendars), extrapolation flagging, and the limits of daily
discharge skill for free-running climate models. \todo{rewrite abstract
once the final objective is formulated; add headline future-projection
result when SSP runs complete}
\end{abstract}

\section{Introduction}
\label{sec:intro}

Scenario-driven water-availability assessments over Indian catchments face a
selection problem that is rarely treated explicitly: which CMIP6 Earth-system
models should feed the hydrological analysis? Existing practice often
inherits ensembles from global assessments \todo{cite examples -- real
literature survey needed}. Three gaps motivate this work:

\begin{enumerate}
\item \textbf{Regional bias structure.} Persistent precipitation biases of
CMIP6 models in the AR6 reference regions covering South Asia
(EAS, WSB, ESB) make raw model output unsuitable for catchment water-balance
studies \todo{cite AR6 WG1 Ch.\,8 / literature -- survey pending}. This
connects to open problem \textbf{C9} (regional climate biases and
explainable downscaling) of the author's research agenda.
\item \textbf{Hydrological (not just meteorological) skill.} A model can
reproduce rainfall statistics yet fail to generate realistic runoff, because
discharge integrates seasonality, soil-moisture dynamics and evaporative
demand. Model selection should therefore include a hydrological test, not
only a forcing test.
\item \textbf{Global-South benchmarking.} Lumped conceptual models remain the
baseline that new (e.g., deep-learning) approaches must beat when local data
are sparse \todo{connect to C21 (Global-South data equity and benchmarking);
survey pending}. Establishing a transparent GR4J baseline over CAMELS-IND
contributes this reference point.
\end{enumerate}

\noindent This manuscript is deliberately structured as a
\emph{methods-and-workflow documentation}: the pipeline is fully
reproducible (Section~\ref{sec:methods}), preliminary results are reported
from the completed historical leg (Section~\ref{sec:results}), and the
future-scenario leg is specified but not yet executed
(Section~\ref{sec:todo}).

\section{Data}
\label{sec:data}

\subsection{CAMELS-IND catchments and observations}
The CAMELS-IND dataset provides daily catchment-mean meteorological forcings
(precipitation, maximum and minimum temperature), observed streamflow
(m\,s$^{-1}$ \todo{unit check: m$^3$/s}), and catchment attributes for 472
catchments in Peninsular India over 1981--2014 \todo{cite the CAMELS-IND
dataset paper -- survey pending}. Forcing availability is complete for all
12,418 days in the study period; streamflow availability varies by gauge.
Catchment area (\texttt{cwc\_area}, km$^2$) and outlet latitude
(\texttt{cwc\_lat}) are taken from the topography attribute file.

\subsection{CMIP6 historical output}
Daily precipitation (\texttt{pr}) and maximum/minimum near-surface
temperature (\texttt{tasmax}, \texttt{tasmin}) were retrieved for 24 CMIP6
models over the Indian domain (1981--2014). Two models (CESM2, IITM-ESM)
provide precipitation only and were excluded; \textbf{22 models} carry all
three variables and form the selection pool. The HadGEM3-GC31-MM 360-day
calendar was handled by mapping its invalid 29/30 February to missing
values. Single first members (\texttt{r1i1p1f1} or closest available) are
used; ensemble spread is not assessed \todo{limitation; survey on large
ensembles pending}.

\section{Methods}
\label{sec:methods}

\subsection{Catchment tiering}
Catchments were tiered by forcing completeness over 1981--2014:
Tier~1 ($\ge$70\% valid days; 180 catchments), Tier~2 ($\ge$30\%; 62), and
Tier~3 ($<$10 years; 230), which is excluded but logged. Every exclusion is
recorded in a machine-readable inventory; nothing is dropped silently.

\subsection{Catchment-mean model series}
Gridded CMIP6 fields were clipped to the Indian boundary and aggregated to
catchment means with an area-weighted average over the intersecting grid
cells, producing one daily series per model$\times$variable$\times$catchment
(66 series files per model).

\subsection{Seasonal empirical quantile mapping (EQM)}
Empirical quantile mapping is fitted per variable, season (DJF, MAM, JJA,
SON) and catchment between model and observed 1981--2014 series:
\begin{itemize}
\item quantile bins require $\ge$300 days; seasons failing the threshold
fall back to the annual fit (331 of 15,972 fits used the fallback);
\item the empirical CDF is tie-collapsed so that constant model values map
consistently;
\item tail policy: precipitation \emph{constant} (clamped to the fitted
range) for the historical evaluation; temperatures \emph{linear}
(extrapolated from the edge bin slopes). For future SSP application the
precipitation policy switches to \emph{linear} to avoid truncating
intensified extremes \todo{cite QM extrapolation literature -- survey
pending};
\item every corrected value carries a flag (\texttt{within\_range} vs.\
extrapolation method); extrapolated values are counted, never hidden.
\end{itemize}

\subsection{Out-of-sample forcing ranking (QED)}
To avoid in-sample optimism, EQM was fitted on 1981--2000 only and applied
to 2001--2014; skill against observed catchment forcings was scored
out-of-sample with KGE, NSE and PBIAS per model$\times$variable, combined
into a composite rank sum. A natural break after rank 14 (non-overlapping
out-of-sample KGE between ranks 1--14 and 15--22 for every variable)
motivated a 14-model shortlist at this stage.

\subsection{GR4J calibration and validation (HYD)}
GR4J \todo{cite Perrin et al.\ 2003} is run with the \texttt{airGR} package
\todo{cite Coron et al.\ 2017} at a daily step. Inputs are precipitation and
potential evapotranspiration computed with the Oudin formula (daily mean
temperature and latitude) \todo{cite Oudin et al.\ 2005}. Observed
discharge is converted to model space, $Q[\mathrm{mm\,d^{-1}}] =
Q[\mathrm{m^3\,s^{-1}}]\times 86.4 / A[\mathrm{km^2}]$.

Per catchment: a 1-year warm-up (1981) precedes the run window (verified:
\texttt{airGR} discards warm-up indices that overlap the run window);
GR4J is calibrated on 1982--2000 with the Michel local-search algorithm
minimising $1-\mathrm{KGE}$ \todo{cite Michel 1991 / airGR design}; and
evaluated \emph{out-of-sample} on 2001--2014 with KGE, NSE and PBIAS
(PBIAS computed manually; \texttt{airGR} provides no such criterion).
Catchments with fewer than 3,650 valid discharge days in the calibration
window are skipped and logged. \todo{Monte-Carlo (LHS) calibration as in
the author's regionalization study is an optional robustness check; survey
pending.}

\subsection{CMIP-driven discharge evaluation}
Each calibrated parameter set is then driven with each of the 22
bias-corrected CMIP6 series (1982--2014, same warm-up), and the simulated
discharge is scored against observed 2001--2014 discharge. Models are ranked
per catchment on KGE, NSE and $|$PBIAS$|$; rank sums across catchments give
a composite \emph{discharge ranking}.

\subsection{Hybrid shortlist synthesis}
The final shortlist combines both rankings (mean of the forcing rank and the
discharge rank). The two worst discharge performers of the forcing-only
top-14 (GISS-E2-1-G: discharge rank 21/22; INM-CM5-0: 20/22) are replaced by
the strongest discharge performers outside it (MPI-ESM1-2-HR: 5/22;
GFDL-ESM4: 9/22), keeping the shortlist at 14 models while retaining the
full CMIP6 ECS range (INM-CM4-8 $\approx$1.8\,K to CanESM5
$\approx$5.7\,K) \todo{confirm ECS values against the literature -- survey
pending}.

\section{Preliminary results (historical leg, completed)}
\label{sec:results}

\begin{table}[h]
\centering\small
\caption{Split-sample performance of calibrated GR4J (192 of 242 Tier-1/2
catchments; 50 skipped for insufficient discharge records, all logged).}
\begin{tabular}{lccc}
\toprule
Window & median KGE & median NSE & share NSE\,$>$\,0.5 \\
\midrule
Calibration 1982--2000 & 0.80 & 0.64 & 0.66 \\
Validation 2001--2014 (out-of-sample) & 0.54 & 0.43 & 0.43 \\
\bottomrule
\end{tabular}
\end{table}

\begin{table}[h]
\centering\small
\caption{Top of the hybrid model shortlist (full table:
\texttt{data\_processed/hyd/model\_shortlist\_final.csv}). QED = forcing
rank, DIS = GR4J discharge rank.}
\begin{tabular}{clccc}
\toprule
Hybrid & Model & QED & DIS & mean volume PBIAS [\%] \\
\midrule
1 & EC-Earth3-Veg-LR & 1 & 3 & $-2.4$ \\
2 & NorESM2-MM       & 7 & 1 & $-4.3$ \\
3 & MIROC6           & 9 & 2 & $+3.9$ \\
4 & MIROC-ES2L       & 3 & 7 & $+17.7$ \\
5 & FGOALS-g3        & 2 & 13 & $+14.1$ \\
6 & HadGEM3-GC31-MM  & 11 & 6 & $+0.4$ \\
7 & EC-Earth3        & 14 & 4 & $+8.9$ \\
\bottomrule
\end{tabular}
\end{table}

\noindent Key observations:

\begin{itemize}
\item Agreement between the forcing and discharge rankings is weak
(Spearman $\rho = 0.14$), confirming that hydrological evaluation adds
information that forcing skill alone does not capture.
\item Absolute daily discharge skill is $\approx$0 for \emph{all} models
(median KGE $-0.05$ to $-0.11$). This is expected: CMIP6 historical runs are
free-running, so their daily weather is not phase-synchronised with
observations; quantile mapping corrects distributions, not timing. Volume
bias (PBIAS) and the \emph{relative} ordering remain the interpretable
signals \todo{support with literature on QM limits for unpaired series --
survey pending}.
\item Models with the smallest discharge volume bias
(NorESM2-MM $-4.3$\%, HadGEM3-GC31-MM $+0.4$\%, MIROC6 $+3.9$\%,
MPI-ESM1-2-HR $+5.3$\%) demonstrate that distribution-corrected series can
reproduce catchment water balances well even where daily timing skill is
bounded.
\item Extrapolation exposure is small but non-negligible: 0.05\% of scored
forcing values and 0.13\% of discharge-evaluation values were flagged as
tail extrapolations; future (warmer) windows will raise this and are the
reason for the \emph{linear} precipitation tail policy.
\end{itemize}

\section{Discussion, limitations and planned work}
\label{sec:todo}

\subsection{Limitations to state honestly}
Single members only; no evaluation of internal-variability spread.
GR4J is calibrated on observed forcings -- parameter transfer to
model climates assumes structural stationarity. Daily timing skill against
free-running models is inherently bounded; seasonal or distributional
hydrological metrics (e.g., monthly flow-duration, low-flow/flood
quantiles) should complement daily KGE in the final analysis. The
discharge-based ranking uses the same evaluation window as the forcing
ranking; a fully independent future window will strengthen the synthesis.

\subsection{Planned extensions (SSP leg)}
For the shortlisted 14 models: download SSP1-2.6, SSP2-4.5 and SSP5-8.5
daily \texttt{pr}/\texttt{tasmax}/\texttt{tasmin} for near-term
(2015--2040), mid-term (2041--2060) and long-term (2081--2100) windows;
apply the historical-fitted EQM with the \emph{linear} precipitation tail
policy; drive calibrated GR4J to produce future daily discharge ensembles
per catchment; derive water-availability and drought-relevant indicators
(e.g., low-flow quantiles, dry-spell statistics) --- connecting to open
problem \textbf{C2} (compound extremes) of the research agenda.

\subsection{Formulation of the final objective}
\todo{This is the open item. Candidate framings, to be resolved after the
literature survey: (i) a methods paper on two-stage GCM selection for
catchment hydrology; (ii) a projection paper on water availability /
drought over Peninsular India under SSPs; (iii) a benchmarking paper
positioning the GR4J baseline for Global-South hydrological ML (C21);
or a combination. See the open-problem list for anchors C9, C21, C2.}

\subsection{Literature survey TODO list}
\todo{GR4J and airGR origins; EQM variants and their behaviour on unpaired
series; CMIP6 Indian-region evaluation studies; CAMELS-IND dataset paper
and prior applications; KGE/NSE interpretation references; ECS values for
the shortlisted models; IAHS HELPING decade positioning (Pangaea-style
data statements).}

\section*{Data and code availability}
The complete workflow is reproducible from R Markdown notebooks
(\texttt{QED\_00--30}: tiering, aggregation, quantile mapping, ranking;
\texttt{HYD\_00--30}: input preparation, GR4J calibration, CMIP-driven
simulation) at \url{https://github.com/anixn/CMIP6_IND_GR4J}. Small result
tables are tracked in the repository; bulk series are regenerated by the
notebooks from raw data.

\begin{thebibliography}{9}
\bibitem{perrin2003} Perrin, C., Michel, C., Andr\'eassian, V. (2003).
Improvement of a parsimonious model for streamflow simulation. Journal of
Hydrology, 279, 275--289.
\bibitem{coron2017} Coron, L., Thirel, G., Guillaumot, S., Le Moine, N.,
Andr\'eassian, V. (2017). airGR: an R package. \todo{verify full citation}
\bibitem{oudin2005} Oudin, L., Hervier, F., Michel, C., Perrin, C.,
Andr\'eassian, V., Anctil, F., Loumagne, C. (2005). Which potential
evapotranspiration input for a lumped rainfall--runoff model? Journal of
Hydrology, 303, 290--306.
\bibitem{eyring2016} Eyring, V., et al. (2016). Overview of the Coupled
Model Intercomparison Project Phase 6 (CMIP6). Geoscientific Model
Development, 9, 1937--1958.
\bibitem{gupta2009} Gupta, H.V., Kling, H., Yilmaz, K.K., Martinez, G.F.
(2009). Decomposition of the mean squared error and NSE performance
criteria. Journal of Hydrology, 377, 80--91.
\bibitem{hargreaves1985} Hargreaves, G.H., Samani, Z.A. (1985). Reference
crop evapotranspiration from temperature. Applied Engineering in
Agriculture, 1, 96--99.
\end{thebibliography}
\noindent\todo{Bibliography is a minimal verified core; complete survey
pending.}

\end{document}
